% Part: computability
% Chapter: computability-theory
% Section: introduction

\documentclass[../../../include/open-logic-section]{subfiles}

\begin{document}

\olfileid{cmp}{thy}{int}
\olsection{مقدمه}

شاخه‌ای از منطق که 
\emph{نظریهٔ محاسبه‌پذیری} نامیده می‌شود به مسائل مربوط به
محاسبه‌پذیری یا محاسبه‌پذیری نسبی توابع و مجموعه‌ها می‌پردازد. اینکه
این حوزه پیش‌تر 
\emph{نظریهٔ بازگشت} نامیده می‌شد نشانه‌ای از نفوذ کلینی است؛ امروزه
هر دو نام رواج دارند.

تابع~$f\colon \Nat \pto \Nat$ را 
\emph{محاسبه‌پذیر جزئی} می‌نامیم اگر در یکی از مدل‌های محاسبه قابل
محاسبه باشد. اگر~$f$ تام باشد، به‌سادگی می‌گوییم~$f$ 
\emph{محاسبه‌پذیر} است. رابطهٔ~$R$ را نیز، اگر تابع مشخصهٔ
محاسبه‌پذیر~$\Char{R}$ داشته باشد، محاسبه‌پذیر می‌نامیم. اگر~$f$ و
$g$ توابع جزئی باشند، $f(x) \fdefined$ می‌نویسیم تا بگوییم~$f$ در~$x$
تعریف شده است؛ یعنی~$x$ در دامنهٔ~$f$ است. $f(x) \fundefined$ معنای
مخالف دارد؛ یعنی~$f$ در~$x$ تعریف نشده است. از
$f(x) \simeq g(x)$ برای بیان این امر استفاده می‌کنیم که یا هر دو
$f(x)$ و~$g(x)$ تعریف‌نشده‌اند، یا هر دو تعریف‌شده و برابرند.

می‌توان نظریهٔ محاسبه‌پذیری را بی‌نیاز از ارجاع به مدل معینی از محاسبه
بررسی کرد. برای این کار نشان می‌دهیم تابع محاسبه‌پذیر جزئیِ جهانی
$\fn{Un}(k,x)$ وجود دارد. این تابع اجازه می‌دهد توابع محاسبه‌پذیر جزئی
را برشماریم. نماد~$\cfind{k}$ را برای تابع محاسبه‌پذیر جزئیِ یک‌موضعی
$k$ام به کار می‌بریم که با
$\cfind{k}(x) \simeq \fn{Un}(k,x)$ تعریف می‌شود. (کلینی برای این
منظور از~$\{k\}$ استفاده می‌کرد، اما این نمادگذاری در سال‌های اخیر
کمتر به کار رفته است.) اندکی کلی‌تر، می‌توانیم توابع محاسبه‌پذیر جزئی
با موضعیت‌های دلخواه را به‌طور یکنواخت برشماریم؛ از~$\cfind{k}[n]$
برای نشان‌دادن تابع بازگشتی جزئیِ~$n$موضعی و~$k$ام استفاده می‌کنیم.

اگر $f(\vec x,y)$ تابعی تام یا جزئی باشد،
$\umin{y}{f(\vec x,y)}$ تابعی از~$\vec x$ است که کوچک‌ترین~$y$ای را
بازمی‌گرداند که $f(\vec x,y)=0$، مشروط بر اینکه همهٔ
$f(\vec x,0)$، \dots، $f(\vec x,y-1)$ تعریف شده باشند. اگر چنین
$y$ای وجود نداشته باشد، $\umin{y}{f(\vec x,y)}$ تعریف‌نشده است. اگر
$R(\vec x,y)$ رابطه باشد، $\umin{y}{R(\vec x,y)}$ را کوچک‌ترین~$y$ای
تعریف می‌کنیم که $R(\vec x,y)$ صادق است؛ به بیان دیگر، کوچک‌ترین~$y$
چنان‌که $1 \tsub \Char{R}(\vec x,y)=0$.

برای نشان‌دادن محاسبه‌پذیری یک تابع، دو راه داریم:
\begin{enumerate}
\item به‌طور صوری: یک ماشین تورینگ یا تابع بازگشتی جزئی را به‌صراحت
  وصف کنیم و نشان دهیم تابع مورد نظر را محاسبه می‌کند؛
\item به‌طور غیرصوری: الگوریتمی را وصف کنیم که آن را محاسبه می‌کند و
  به تز چرچ استناد کنیم.
\end{enumerate}
مرز دقیقی میان این دو نیست. وصف تفصیلی یک الگوریتم باید اطلاعات کافی
فراهم کند تا نسبتاً روشن باشد چگونه می‌توان، دست‌کم در اصل، ماشین
تورینگ یا دنبالهٔ مناسبی از تعریف‌های بازگشتی جزئی طراحی کرد. بعید است
تعریف‌های کاملاً صوری روشنگر باشند؛ ازاین‌رو می‌کوشیم حد میانهٔ مناسبی
میان این دو رویکرد بیابیم. خلاصه، می‌کوشیم اثبات‌هایی شهودی اما دقیق
به دست دهیم که نشان دهند تعریف‌های صریح را می‌توان ساخت.

\end{document}
